% !TEX root = ../main.tex
\clearpage
\phantomsection
\section*{10 \textbf{학위논문 개요}}
\label{sec:dissertationthesis-outline}
\addcontentsline{toc}{section}{10 학위논문 개요}

학위논문은 다음 여섯 장으로 구성된다:


\begin{itemize}
\item[\textendash] 제1장: 서론 및 연구 문제. 본 장은 계획서를 학위논문 서론으로 정제하여 Arbitrum One의 DeFi 레버리지 거래 맥락, 온체인 평판 점수 문제, 연구 질문, 목표, 범위, 의의, 한계를 제시한다.
\item[\textendash] 제2장: 문헌 검토 및 이론적 틀. 본 장은 계획서 문헌 검토를 그래프 기반 평판 점수 산출, DeFi 신용 리스크 예측, ERC-20 승인 의미론, PageRank 변형, 평판 구조·GMX 거래 성공 정렬·계산 효율성에 대한 이론적 토대에 대한 포괄적 분석으로 확장한다.
\item[\textendash] 제3장: 연구 방법론. 본 장은 철학적 입장, 연구 설계, 데이터 소스, 표본 추출 논리, 특성 공학, 그래프 구성 절차, 기준선 재현, 통계 검정, 견고성 검사, 윤리 심의, 재현성 통제를 명시한다.
\item[\textendash] 제4장: 구현 및 실증 결과. 본 장은 데이터 파이프라인, EndorseRank 구현, 기준선 구현, 데이터셋 특성, 그래프 통계, GMX 거래 성공 정렬 결과, 순위 발산 결과, 계산 벤치마크 결과를 보고한다.
\item[\textendash] 제5장: 논의. 본 장은 각 연구 질문 및 목표에 대해 결과를 해석하고, DeFi 리스크 평가에 대한 이론적·실무적 함의를 설명하며, 타당성 위협을 평가하고, 결과가 주장된 기여를 어떻게 지지하거나 제한하는지 논의한다.
\item[\textendash] 제6장: 결론 및 향후 연구. 본 장은 연구를 종합하고, 지식에 대한 독창적 기여를 명시하며, 연구 질문을 재검토하고, 한계를 요약하며, 교차 프로토콜 검증, 적대적 조작, 더 넓은 온체인 신뢰 신호에 대한 향후 연구를 제안한다.
\end{itemize}
\phantomsection
